\documentclass[11pt,a4paper]{article}
\usepackage{fontspec}\usepackage[french]{babel}
\usepackage[margin=1.8cm]{geometry}\usepackage{tabularx,colortbl,xcolor,array}
\setmainfont{PlusJakartaSans-Medium.ttf}[Path=fonts/,BoldFont=PlusJakartaSans-Bold.ttf]
\pagestyle{empty}\setlength{\parindent}{0pt}
\begin{document}
{\LARGE\bfseries Espagnol — Terminale A\par}\smallskip
{\large Répartition des 29 leçons par module et par trimestre\par}\medskip
Cours « signés OnBuch+ », conformes à la fiche de progression harmonisée du MINESEC (année scolaire 2026-2027).
Chaque leçon est un PDF autonome (2544 pages au total pour les 29 leçons).\par\medskip
\textbf{Contenu type d'une leçon :} page de garde, sommaire, objectifs, prérequis, sections de cours (textes en espagnol avec glossaire,
règles de grammaire, tableaux, exemples traduits), activité d'intégration, méthodes et exercices résolus, exercices (connaissances, entraînement,
préparation à l'examen), corrigés, fiche bilan. Types de leçon : \emph{texte} (lecture et commentaire, lexique), \emph{traduction} (grammaire,
thème et version), \emph{mixte}.\par\medskip
\textbf{Les 5 modules :} 1. Vie familiale et intégration sociale (leçons 1--6) ; 2. Environnement, santé et bien-être (7--11) ;
3. Vie citoyenne et ouverture au monde (12--15) ; 4. Activités économiques et monde du travail (16--21) ;
5. Médias, communication et culture de la paix (22--29).\par\bigskip
\section*{Module 1 — Vie familiale et intégration sociale}
\textit{6 leçons, leçons 1 à 6, 530 pages.}\par\medskip
\begin{tabularx}{\linewidth}{|c|X|l|l|c|}\hline
\rowcolor{gray!20} N° & Titre de la leçon & Trimestre & Période & Pages \\ \hline
1 & Lectura y comentario : discriminations et violences contre les femmes (la familia hispánica) \newline {\small\emph{Type : texte — fichier : Lecon01-discriminations-violences-femmes.pdf}} & 1er trimestre & 07/09 – 16/10/2026 & 104 \\ \hline
2 & Diálogo familiar, debate sobre la violencia de género et traduction : SER y ESTAR \newline {\small\emph{Type : mixte — fichier : Lecon02-dialogo-familiar-ser-estar.pdf}} & 1er trimestre & 07/09 – 16/10/2026 & 83 \\ \hline
3 & Lectura y comentario : derechos de las mujeres y de las minorías \newline {\small\emph{Type : texte — fichier : Lecon03-derechos-mujeres-minorias.pdf}} & 1er trimestre & 07/09 – 16/10/2026 & 82 \\ \hline
4 & Traduction de « C'est… que », « C'est… qui », « Ne… que » et de l'emphase ; redacción : describir a su familia (80 palabras) \newline {\small\emph{Type : traduction — fichier : Lecon04-traduccion-enfasis-descripcion-familia.pdf}} & 1er trimestre & 07/09 – 16/10/2026 & 86 \\ \hline
5 & Lectura y comentario : la familia y sus problemas (viudez, herencia) ; juego de rol \newline {\small\emph{Type : texte — fichier : Lecon05-familia-problemas-viudez-herencia.pdf}} & 1er trimestre & 07/09 – 16/10/2026 & 90 \\ \hline
6 & Traductions : l'ordre, l'interdiction, « chez », « devenir » ; debates : educación de los niños y NTIC \newline {\small\emph{Type : mixte — fichier : Lecon06-orden-prohibicion-chez-devenir.pdf}} & 1er trimestre & 07/09 – 16/10/2026 & 85 \\ \hline
\end{tabularx}\bigskip
\section*{Module 2 — Environnement, santé et bien-être}
\textit{5 leçons, leçons 7 à 11, 439 pages.}\par\medskip
\begin{tabularx}{\linewidth}{|c|X|l|l|c|}\hline
\rowcolor{gray!20} N° & Titre de la leçon & Trimestre & Période & Pages \\ \hline
7 & Lectura y comentario : « El medioambiente y nosotros » (residuos, recogida, reciclaje) \newline {\small\emph{Type : texte — fichier : Lecon07-medioambiente-residuos-reciclaje.pdf}} & 1er trimestre & 19/10 – 30/10/2026 & 78 \\ \hline
8 & Traduction : la phrase impérative avec ou sans pronoms compléments ; emisión de salud / podcast ecológico \newline {\small\emph{Type : mixte — fichier : Lecon08-imperativo-pronombres-podcast.pdf}} & 1er trimestre & 19/10 – 30/10/2026 & 83 \\ \hline
9 & Lectura y comentario : contaminación y medidas de preservación ; creación de carteles de sensibilización \newline {\small\emph{Type : texte — fichier : Lecon09-contaminacion-carteles-sensibilizacion.pdf}} & 1er trimestre & 02/11 – 13/11/2026 & 91 \\ \hline
10 & Traduction : les locutions adverbiales ; redacción : ensayo argumentativo sobre el medioambiente \newline {\small\emph{Type : mixte — fichier : Lecon10-locuciones-adverbiales-ensayo-ambiente.pdf}} & 1er trimestre & 02/11 – 13/11/2026 & 101 \\ \hline
11 & Lectura y comentario de poema ; figuras retóricas ; debate : alimentación sana vs comida basura \newline {\small\emph{Type : mixte — fichier : Lecon11-poema-alimentacion-sana-comida-basura.pdf}} & 1er trimestre & 02/11 – 13/11/2026 & 86 \\ \hline
\end{tabularx}\bigskip
\section*{Module 3 — Vie citoyenne et ouverture au monde}
\textit{4 leçons, leçons 12 à 15, 341 pages.}\par\medskip
\begin{tabularx}{\linewidth}{|c|X|l|l|c|}\hline
\rowcolor{gray!20} N° & Titre de la leçon & Trimestre & Période & Pages \\ \hline
12 & Traduction de quelques vers d'un poème ; POR y PARA ; debate guiado : « Ciudadanos del mundo » \newline {\small\emph{Type : mixte — fichier : Lecon12-por-y-para-poema-ciudadanos-del-mundo.pdf}} & 2e trimestre & 16/11 – 18/12/2026 & 100 \\ \hline
13 & Lectura y comentario : relaciones internacionales, relaciones Norte-Sur, sistemas políticos, derechos y deberes \newline {\small\emph{Type : texte — fichier : Lecon13-relaciones-internacionales-norte-sur.pdf}} & 2e trimestre & 16/11 – 18/12/2026 & 72 \\ \hline
14 & Traduction : verbes défectifs, phrase restrictive, discours indirect \newline {\small\emph{Type : traduction — fichier : Lecon14-verbos-defectivos-restrictiva-estilo-indirecto.pdf}} & 2e trimestre & 16/11 – 18/12/2026 & 78 \\ \hline
15 & Traduction : expression de l'accord et du désaccord \newline {\small\emph{Type : traduction — fichier : Lecon15-acuerdo-y-desacuerdo.pdf}} & 2e trimestre & 16/11 – 18/12/2026 & 91 \\ \hline
\end{tabularx}\bigskip
\section*{Module 4 — Activités économiques et monde du travail}
\textit{6 leçons, leçons 16 à 21, 536 pages.}\par\medskip
\begin{tabularx}{\linewidth}{|c|X|l|l|c|}\hline
\rowcolor{gray!20} N° & Titre de la leçon & Trimestre & Période & Pages \\ \hline
16 & Lectura y comentario : « El mundo laboral hoy » ; ayudas, subvenciones, acuerdos bilaterales y multilaterales \newline {\small\emph{Type : texte — fichier : Lecon16-mundo-laboral-ayudas-desarrollo.pdf}} & 2e trimestre & 04/01 – 12/02/2027 & 96 \\ \hline
17 & Traduction : comparatifs et superlatifs ; redacción : el currículum vítae \newline {\small\emph{Type : mixte — fichier : Lecon17-comparativos-superlativos-curriculum.pdf}} & 2e trimestre & 04/01 – 12/02/2027 & 94 \\ \hline
18 & L'emphase ; redacción : carta de motivación (120 palabras) \newline {\small\emph{Type : mixte — fichier : Lecon18-enfasis-carta-de-motivacion.pdf}} & 2e trimestre & 04/01 – 12/02/2027 & 70 \\ \hline
19 & Redacción : carta abierta / petición ; presentación oral de un país hispanohablante \newline {\small\emph{Type : mixte — fichier : Lecon19-carta-abierta-peticion-pais-hispanohablante.pdf}} & 2e trimestre & 04/01 – 12/02/2027 & 86 \\ \hline
20 & Lectura y comentario : oficios, mundo del trabajo, derechos y deberes de los trabajadores \newline {\small\emph{Type : texte — fichier : Lecon20-oficios-derechos-deberes-trabajadores.pdf}} & 2e trimestre & 04/01 – 12/02/2027 & 82 \\ \hline
21 & Traduction : phrase causale, finalité, futur dans la subordonnée, emphase (révision) \newline {\small\emph{Type : traduction — fichier : Lecon21-causal-finalidad-futuro-subordinada.pdf}} & 2e trimestre & 04/01 – 12/02/2027 & 108 \\ \hline
\end{tabularx}\bigskip
\section*{Module 5 — Médias, communication et culture de la paix}
\textit{8 leçons, leçons 22 à 29, 698 pages.}\par\medskip
\begin{tabularx}{\linewidth}{|c|X|l|l|c|}\hline
\rowcolor{gray!20} N° & Titre de la leçon & Trimestre & Période & Pages \\ \hline
22 & Lectura y comentario : « Los medios de comunicación » ; textes argumentatifs / injonctifs ; cultura de la paz \newline {\small\emph{Type : texte — fichier : Lecon22-medios-de-comunicacion-cultura-de-paz.pdf}} & 3e trimestre & 13/02 – 07/05/2027 & 93 \\ \hline
23 & Traduction : les traductions du « ON » ; les adverbes d'intensité (très, trop, beaucoup, assez) \newline {\small\emph{Type : traduction — fichier : Lecon23-traducciones-on-adverbios-intensidad.pdf}} & 3e trimestre & 13/02 – 07/05/2027 & 88 \\ \hline
24 & Lectura y comentario : el cuarto poder y la promoción de la cultura de paz \newline {\small\emph{Type : texte — fichier : Lecon24-cuarto-poder-cultura-de-paz.pdf}} & 3e trimestre & 13/02 – 07/05/2027 & 92 \\ \hline
25 & Traduction : « avoir beau » ; « demander » ; « depuis » \newline {\small\emph{Type : traduction — fichier : Lecon25-aunque-pedir-desde-hace.pdf}} & 3e trimestre & 13/02 – 07/05/2027 & 71 \\ \hline
26 & Lectura y comentario : libertad de expresión y delitos de prensa ; análisis de publicidad y fake news \newline {\small\emph{Type : texte — fichier : Lecon26-libertad-expresion-delitos-prensa-fake-news.pdf}} & 3e trimestre & 13/02 – 07/05/2027 & 97 \\ \hline
27 & Traduction : la concordance des temps ; le style direct et indirect \newline {\small\emph{Type : traduction — fichier : Lecon27-concordancia-tiempos-estilo-directo-indirecto.pdf}} & 3e trimestre & 13/02 – 07/05/2027 & 85 \\ \hline
28 & Lectura y comentario : Internet y las NTIC ; traduction : expression de la préférence et de l'habitude \newline {\small\emph{Type : mixte — fichier : Lecon28-internet-ntic-preferencia-habito.pdf}} & 3e trimestre & 13/02 – 07/05/2027 & 89 \\ \hline
29 & Redacción : artículo de prensa (120 palabras) ; traduction : la condition réalisable et irréalisable \newline {\small\emph{Type : mixte — fichier : Lecon29-articulo-de-prensa-condicion-realizable-irrealizable.pdf}} & 3e trimestre & 13/02 – 07/05/2027 & 83 \\ \hline
\end{tabularx}\bigskip
\end{document}
